Dans un cyclotron, des protons ($m_p = \qty{1.673e-27}{kg}$, $e = \qty{1.602e-19}{C}$) décrivent des demi-cercles de plus en plus grands dans un champ magnétique uniforme de \qty{\v{B}}{T}. Ils sortent de l'appareil lorsque le rayon de leur trajectoire atteint \qty{\v{R}}{m}.

\begin{enumerate}[label=\alph*)]
    \item Quelle est la vitesse des protons à la sortie ?
    \item Quelle est leur énergie cinétique à la sortie, en mégaélectronvolts ?
    \item À quelle fréquence la tension alternative entre les deux électrodes en D doit-elle osciller ?
\end{enumerate}
